\subsection{Motivation}
\label{sec:motivation}

\IEEEPARstart{C}{areer} decision-making is increasingly shaped by cross-profession and cross-sector mobility into labor markets that are both complex and difficult to navigate.
This thesis is motivated by that tension in the software labor market, and by the resulting need for objective, evidence-based career intelligence rather than generic advice.

A first source of friction is informational complexity.
Occupational labels in software are numerous, and vacancies that share a similar title often request substantially different skill bundles.
Prior analyses of early-career information systems advertisements already documented heterogeneous clusters of technical and non-technical requirements within ostensibly related roles~\cite{Kennan2008}.
More recent work shows that Science, Technology, Engineering, and Mathematics (STEM) vacancies change especially quickly: technological change introduces new job skills and renders older ones obsolete, which helps explain perceived STEM shortages together with rapid skill churn over workers' careers~\cite{DemingNoray2018}.
Complementary survey evidence confirms that online job advertisements are unstructured text and that skill identification from such postings remains a non-trivial research problem, with substantial variation in skill bases, extraction methods, and granularity~\cite{Khaouja2021}.
In parallel, labor-market returns have increasingly favored combinations of technical and social skills, further widening the gap between a job title and the portfolio of competencies that employers actually screen for~\cite{Deming2017}.
Taken together, these findings imply that entrants and career changers face an opaque signal: the market is not a single stable ``programmer'' profile, but a moving distribution of role-specific demands.

A second source of friction is difficulty at the point of entry for early-career professionals.
Comparative survey evidence indicates that faculty in information technology (IT)-related programmes and IT managers in industry do not fully agree on which knowledge and skills matter most for entry-level workers, which complicates preparation for the first job~\cite{Aasheim2009}.
Against heterogeneous vacancy requirements and rapid STEM skill change~\cite{Kennan2008,DemingNoray2018}, newcomers therefore face a high informational and strategic cost when deciding where to invest scarce learning time.
Task-based models of technical change emphasize that technology does not eliminate occupations wholesale; it reallocates tasks between capital and labor and reshapes the skill content of work~\cite{AcemogluAutor2010,AcemogluRestrepo2019}.
Recent administrative evidence is consistent with that reallocation biting hardest at the entry point: after the diffusion of generative artificial intelligence, early-career workers (ages 22--25) in the most AI-exposed occupations---including software development---experienced a roughly 16\% relative employment decline compared with less exposed occupations, even after controlling for firm-level shocks, while employment for more experienced workers in the same fields remained more stable~\cite{BrynjolfssonChandarChen2025}.
Complementary vacancy data for software roles likewise show a post-ChatGPT decline of about 14--15\% in junior postings relative to senior ones~\cite{ModestinoWestbyCheng2026}.
Together, these patterns suggest that the current conjuncture of programming assistants does not merely change tools; it tightens the labor-market doorway that early-career entrants must pass through.

Despite that difficulty, people continue to enter programming pathways.
This apparent paradox can be interpreted through the economic logic of rebound effects---classically associated with Jevons' observation that efficiency improvements may increase rather than reduce overall resource use~\cite{vandenBergh2010}---applied here by analogy to software production.
If tools that raise the productivity of producing software make software cheaper relative to alternatives, demand for software-intensive products and services expands, and new complementary tasks emerge.
In the labor framework of Acemoglu and Restrepo, such expansion corresponds to a reinstatement channel that can sustain or grow demand for human work even as specific tasks are automated~\cite{AcemogluRestrepo2019}.
On a medium-term horizon, therefore, software-related competence remains a rational destination for many knowledge-work transitions, even while long-horizon automation may eventually reach further into cognitive tasks if reinstatement weakens.

These two motives are complementary rather than contradictory.
Structural demand for software-related skills continues to pull workers toward the field, while skill heterogeneity, rapid obsolescence, and shifting task content make the market costly to interpret---especially for early-career entrants and those migrating from other sectors.
